\section{The clocked component-radius threshold}\label{sec:intro54}
A finite stochastic register can carry infinitely many probability distributions. Thus an infinite numerical orbit is not, by itself, an obstruction to finite positive realization. A further difficulty arises when the realization may be redesigned at every horizon and its transitions may depend on a free external clock. The obstruction established here survives both freedoms. It is measured by the radius of an observable image inside its allowed output set.

For the clocked theorem in this section, let $H$ be a compact Hausdorff group with finitely many connected components. Put $K=H^\circ$ and $C=H/K$. Then $K$ is an open normal subgroup and $C$ is a finite group. Fix a finite alphabet $\mathcal A$ with elements $u_a\in H$ whose generated group is dense in $H$. Our convention is
\[
 u_{a_1\cdots a_n}=u_{a_n}\cdots u_{a_1}.
\]
The identity $e$ need not be an alphabet element. Instead assume the following executable return condition:
\begin{equation}\label{eq:returns54}
 \text{there is an integer }g\ge0\text{ such that, for every }n\ge g,
 \text{ some word of length }n\text{ has product }e.
\end{equation}
An identity letter gives $g=0$. The condition concerns physical products; the stochastic rows on a return word may perform arbitrary processing.

For nonempty finite seed and query sets $\mathcal S,\mathcal J$, let
\[
 F_{s,j}:H\longrightarrow Y_j
\]
be continuous, where $Y_j$ is a nonempty compact convex subset of a finite-dimensional real normed space $(V_j,\norm{\cdot}_j)$. A query is selected only after the command word. A prescribed joint law of finitely many outputs is treated as one query, not as a collection of marginal constraints. In particular, a categorical interface has
\[
 Y_j=\Delta_{M_j}=\{p\in\R^{M_j}:p_i\ge0,\ \textstyle\sum_i p_i=1\},
 \qquad \norm{x}_j=\tfrac12\sum_i|x_i|.
\]
Then $\norm{p-q}_j$ is total variation. Binary means correspond to $Y=[-1,1]$ with $\norm{x}=|x|/2$.

\begin{definition}[Clocked realization]\label{def:machine54}
At horizon $N$ choose finite nonempty registers $S_t$, $0\le t\le N$, initialization rows $E_s$, row-stochastic command matrices $T_{t,a}:S_{t-1}\to S_t$, and terminal decoder maps $D_j:S_N\to Y_j$. The realized output vector on $w=a_1\cdots a_N$ is
\[
 E_sT_{1,a_1}\cdots T_{N,a_N}D_j.
\]
Its error is the maximum, over all seeds, length-$N$ words, and queries, of its distance from $F_{s,j}(u_w)$. The command width is $\max_t|S_t|$, counting all declared labels, including unused padding labels. Write $W_{N,\varepsilon}$ for its least possible value at error at most $\varepsilon$.
\end{definition}
All retained private information, including randomness kept for future use, is charged to the current label. The rows may depend on $t$ and $N$, but not on an unrecorded input history or on the future query. The external clock and horizon, arbitrary real table descriptions, their construction and lookup, real arithmetic, and exact sampling are uncharged. This is a numerical realization invariant, not uniform workspace or program size. For a finite categorical interface, charging the terminal answer register changes the peak to $\max(W_{N,\varepsilon},\max_j M_j)$. Thus all boundedness and growth results are unchanged; the convention that also charges a sampled binary-answer register has a minimum of two labels. A decoder vector by itself is not an additional sampled register.

For $A\subset Y_j$ define its constrained radius by
\[
 \rad_{Y_j}(A)=\min_{y\in Y_j}\sup_{x\in A}\norm{x-y}_j.
\]
The minimum is attained when $A$ is compact. The critical error is
\begin{equation}\label{eq:radius54}
 \epsc=\max_{s,j,c\in C}\rad_{Y_j}\bigl(F_{s,j}(c)\bigr).
\end{equation}
Here a component is identified with its coset in $H$; its image is compact by continuity. The requirement that the center belong to $Y_j$ is essential for categorical outputs.

\begin{theorem}[Intrinsic radius classification]\label{thm:radius54}
Under \eqref{eq:returns54}, for every $\varepsilon\ge0$ the following are equivalent:
\begin{enumerate}
\item $\sup_N W_{N,\varepsilon}<\infty$;
\item one stationary finite-state realization with permutation command updates has error at most $\varepsilon$ at every horizon;
\item $\varepsilon\ge\epsc$.
\end{enumerate}
At the boundary, at most $|\mathcal S||C|$ command labels suffice. If $\varepsilon<\epsc$, then for every integer $k\ge1$ there is an integer $L_k<\infty$ such that every error-$\varepsilon$ realization satisfies
\begin{equation}\label{eq:occupation54}
 \#\{t\in\{1,\ldots,N\}:|S_t|\le k\}\le L_k.
\end{equation}
All intervening widths are unrestricted. In particular, $W_{N,\varepsilon}\to\infty$.
\end{theorem}
Neither a positive lower bound on state masses nor a reachable-rank, observability, or conditioning assumption occurs in the theorem. The upper realization is deterministic until its terminal decoder. It need not have the minimum width among all stochastic realizations. Theorem~\ref{thm:stationaryminimum55}, rather than a deterministic component partition, characterizes that stationary minimum.

For binary means, \eqref{eq:radius54} reduces to one quarter of the largest componentwise oscillation. For a genuine categorical law it is generally not half its diameter; an analytic three-outcome example below has critical total-variation error $2/3$, whereas every individual event has binary threshold $1/2$. The proof uses finitely many localized vector tests rather than a separate midpoint for each coordinate.

The finite-component assumption does not assert that $H$ is a Lie group. For example, take $H=\prod_{n\ge1}\R/\Z$, choose an element whose every finite coordinate tuple together with $1$ is rationally independent, and include this element and $e$ as commands. The positive orbit is dense and $H$ is connected but has no faithful finite-dimensional representation. The theorem applies to every finite continuous interface on this group. With infinitely many components the component construction need not be finite; Section~\ref{sec:profinite55} gives the finite-observable-quotient criterion and the strict-side clocked theorem in that setting.

There is also a quantitative counterpart. For planar rotations with a dual badly approximable $s$-tuple of angles and signal amplitude $0<\rho<1$, Theorem~\ref{thm:matched54} proves
\[
 W_{N,\varepsilon}=\Theta\bigl(N^{s/(2s+1)}\bigr)
 \quad\text{for every fixed }0\le\varepsilon<\rho/2.
\]
The upper realization is exact even when $\varepsilon>0$. The new lower estimate reaches the entire subcritical interval, by combining harmonic localization with the arithmetic orbit-packing mechanism. This quantitative theorem and the radius classification have distinct hypotheses: algebraic matrix entries alone do not imply the badly approximable angle condition.

\subsection{Executable tests}
Let $\mathcal P=\{u_w:w\in\mathcal A^*\}$. We use standard Haar averaging and the density of matrix coefficients of finite-dimensional unitary representations, in realified form \cite{Folland}. No quantitative mixing assumption for a preassigned random walk is imposed.

\begin{lemma}[Positive words and common lengths]\label{lem:positive54}
The semigroup $\mathcal P$ is dense in $H$, and $\mathcal P\cap K$ is dense in $K$. Every finite subset of $\mathcal P\cap K$ has executable representatives of a common positive length.
\end{lemma}
\begin{proof}
For an element $u$ of a compact group, positive powers return arbitrarily close to $e$ with exponent at least two. Indeed cover the compact closure of its powers by finitely many translates of a sufficiently small neighborhood. Among sufficiently many powers, three lie in one member of the cover; the first and last exponents differ by at least two, and their quotient is arbitrarily close to $e$. This argument also covers finite-order elements by taking a sufficiently large multiple of the order. If $u^r$ approaches $e$ with $r\ge2$, then $u^{r-1}=u^ru^{-1}$ approaches $u^{-1}$ and is still a positive power. Thus $u^{-1}$ is in the closure of positive powers. The closed semigroup $\overline{\mathcal P}$ therefore contains the group generated by the letters, and hence equals $H$. Since $H$ has finitely many components, $K$ is open, and a relatively open subset of $K$ is open in $H$; density gives the second assertion. For finitely many representative lengths $\ell_i$, choose $B\ge\max_i\ell_i+g$, $B\ge1$, and append return words of lengths $B-\ell_i$. The products are unchanged.
\end{proof}

\begin{lemma}[Finite executable gap]\label{lem:gap54}
Let $R:K\to O(E)$ be a continuous representation on a nonzero finite-dimensional Euclidean space, with no nonzero fixed vector. For every $k\ge1$ there are $B\ge1$, a probability law $\nu$ on finitely many executable length-$B$ words with products in $K$, and $d\in(0,1)$ such that
\begin{equation}\label{eq:gap54}
 \E_{w\sim\nu}\max_{1\le i\le k}\ip{v_i}{R(u_w)u}\le1-d
\end{equation}
for all unit vectors $u,v_1,\ldots,v_k\in E$.
\end{lemma}
\begin{proof}
On the compact product of unit spheres consider the minimum
\[
 \gamma=\min_{u,v_1,\ldots,v_k}\int_K
 \left(1-\max_i\ip{v_i}{R(h)u}\right)\,d\mu(h).
\]
Repeated centers are permitted. The integrand is nonnegative. If a minimizing tuple had integral zero, full support of Haar measure would give $R(h)u\in\{v_1,\ldots,v_k\}$ for every $h$. The orbit of $u$ is connected, hence a singleton, contradicting the fixed-vector hypothesis. Thus $\gamma>0$.

Partition the compact image $R(K)$ into finitely many measurable sets of small diameter in the Euclidean operator norm, and replace their Haar masses by atoms at nearby executable products in $K$. Cells of zero Haar mass can be discarded. The function $A\mapsto\max_i\ip{v_i}{Au}$ is uniformly $1$-Lipschitz in that norm. An approximation radius less than $\gamma/2$ therefore preserves a positive defect uniformly in the unit vectors. Lemma~\ref{lem:positive54} gives a common length. Decreasing the defect to a number in $(0,1)$ proves the assertion.
\end{proof}

\begin{lemma}[Conditional-centroid contraction]\label{lem:centroid54}
Let $Z\in E$ be a bounded random physical vector and $S$ the current register. Apply an independent word $w\sim\nu$ from Lemma~\ref{lem:gap54}, and let the actual block kernel generate a terminal state $S'$ having at most $k$ labels. Then
\begin{equation}\label{eq:centroid54}
 \E\norm{\E[R(u_w)Z\mid S']}\le(1-d)\E\norm{\E[Z\mid S]}.
\end{equation}
A fixed intervening return word cannot increase the same conditional norm, regardless of its intermediate or terminal widths.
\end{lemma}
\begin{proof}
Put $z_s=\E[Z\mid S=s]$ and write $T_w(s,i)$ for the actual composite row of the block. Since this row depends on the past only through $S$,
\[
 \Pr(S'=i)\E[R(u_w)Z\mid S'=i]
 =\sum_s\Pr(S=s)\E_w T_w(s,i)R(u_w)z_s.
\]
Choose $v_i$ in the direction of each nonzero terminal centroid, arbitrarily otherwise. Pair, sum in $i$, and use stochasticity to bound the sum by
\[
 \sum_s\Pr(S=s)\E_w\max_i\ip{v_i}{R(u_w)z_s}
 \le(1-d)\sum_s\Pr(S=s)\norm{z_s}.
\]
Zero-probability states contribute zero. No restriction was placed on registers inside the block. On a fixed return word the physical vector is unchanged at its endpoints; the new state is a stochastic image of the old one, so conditional Jensen gives nonincrease. Independence is used only at complete block boundaries.
\end{proof}
